% !TeX root = ../thesis.tex
\section{Conclusion}
\label{sec:conclusion}

This study demonstrated that the final accuracy of the deployed Text-to-SQL system 
did not fully describe how it answered questions and reached its results.
Correct results could be produced after an unsuccessful query execution, 
and revised queries could still yield the wrong answer.
By capturing query execution traces, database feedback and result correctness, 
it was possible to characterise this discrepancy.
The comparison of the models across repeated runs showed 
behaviour that their final scores alone did not explain.

The main contribution was an evaluation methodology that leveraged 
the captured traces to help explain how the system arrived at both correct and incorrect results. 
It allowed the examination of individual failures alongside patterns across questions and runs. 
This could support the debugging process by detecting repeated errors, 
unsuccessful query revisions and unfulfilled requirements. 
If extended to a broader set of questions and databases relevant to the deployment, 
the methodology could help with model selection by accounting for answer accuracy, 
query error recovery and consistency.

Further evaluation using additional datasets and repeated runs would be necessary to confirm these uses. 
Nevertheless, the gathered results showed that the captured traces provided insights into how 
the system produced its answers beyond the accuracy score alone—in particular, 
whether query execution recovered from an error or whether additional analysis was required.